\subsubsection{The exact pure-point-complex boundary}

The preceding obstruction extends beyond a single cohomological
parity. The relevant constraint is the cohomological shift in the
definition of purity.

\begin{finding}[\fid{F57}, pure-complex refinement]
\label{f:purepointLaurentcriterion}
Let \(q\) be a prime power, \(\ell\nmid q\), and choose
\(\lambda=\sqrt q>0\) under a specified complex embedding,
together with its realization in a finite \(\ell\)-adic
coefficient field \(E\).
For
\[
 f(v)=\sum_e c_ev^e\in\mathbb Z[v,v^{-1}],\qquad w\in\mathbb Z,
\]
the sequence \(f(q^{n/2})\) is the Frobenius Euler trace of a
bounded Weil complex on \(\operatorname{Spec}\mathbb F_q\)
pure of the single weight \(w\) if and only if
\begin{equation}\label{eq:purepointLaurentsign}
                 (-1)^w f(-v)\in\mathbb N[v,v^{-1}].
\end{equation}
Necessity holds even if equality is required only for all
sufficiently large \(n\); sufficiency gives equality for every
positive \(n\). The inverse convention \(v=q^{-n/2}\) has the
same sign test.

Consequently the actual word trace
\eqref{eq:rankfournegativeLaurent} cannot be the Euler trace
of any bounded pure point complex of any single weight.
The same holds for the actual degeneration coefficient
\[
                   c_0(v)=v^2+v^{-2}-2.
\]
Both conclusions permit mixed complexes. They do not prohibit
pure coefficients on a nonproper source whose compact
pushforward to the point is mixed.
\end{finding}

\begin{proof}
BBD, \S5.1.5, pp.~125--126, defines punctual purity using
integral weights and all complex conjugates.
The smooth/lisse clause of \S5.1.8, pp.~126--127, says that
a complex \(K\) is pure of weight \(w\) precisely when
each \(H^iK\) is punctually pure of weight \(w+i\)
\cite{BBD}. On the point every constructible sheaf is lisse.

Write \(m_{i,\alpha}\) for the algebraic multiplicity of
the nonzero Frobenius eigenvalue \(\alpha\) on \(H^iK\).
Jordan blocks do not affect power traces, so
\[
 \operatorname{Tr}_{\mathrm{Eul}}(F^n;K)
      =\sum_{i,\alpha}(-1)^i m_{i,\alpha}\alpha^n.
\]
Group equal roots and compare with
\(\sum_e c_e\lambda^{en}\).
Finite Vandermonde independence, including its
eventual-equality version in \fid{F02}, identifies the
coefficient at each root \(\lambda^e\).
Purity permits this root only in degree
\[
                            i=e-w,
\]
since its absolute value is \(q^{e/2}\). Hence
\[
                    c_e=(-1)^{e-w}m_{e-w,\lambda^e}.
\]
This proves \eqref{eq:purepointLaurentsign}.
It also rules out hidden cancellation: any fixed root
has a fixed absolute value and therefore can occur in at
most one degree of a complex pure of one weight.
There are no opposite-parity contributions at that root
which could cancel. Arithmetic Frobenius need not be
semisimple for this argument.

Conversely, put \(m_e=(-1)^{e-w}c_e\ge0\).
Let \(L_e\) be the one-dimensional \(E\)-Weil line with
geometric Frobenius \(\lambda^e\). Every algebraic conjugate
has absolute value \(q^{e/2}\), so \(L_e\) is pure of
weight \(e\). Then
\begin{equation}\label{eq:purepointLaurentconstruction}
           K=\bigoplus_e L_e[-(e-w)]^{\oplus m_e}
\end{equation}
has \(H^iK\) pure of weight \(w+i\), and its Euler trace
is \(f(q^{n/2})\). The line can also be chosen continuous
as a Galois representation: \(\lambda^e\) is an
\(\ell\)-adic unit, so the cyclic image has compact
closure and extends to the procyclic Galois group.
The point-representation convention is BBD,
\S5.1.11, pp.~127--128.
For even \(e=2k\), one may use
\(L_{2k}=\mathbb Q_\ell(-k)\).
For odd \(e\) this is an arithmetic coefficient line,
not a claim of a geometric half-Tate motive.

For the inverse variable convention the possible degree
is \(-e-w\), which has the same parity as \(e-w\).
All exponents of \eqref{eq:rankfournegativeLaurent}
are even, so the sign test would require every coefficient
to have the same sign \((-1)^w\).
Its coefficients \(-1\) at the two outer exponents and
\(668\) at exponent zero exclude both parities.
The coefficients \(1,-2,1\) of \(c_0\) do the same.
An overall shift or a common Tate twist does not repair
these mixed signs.
\end{proof}

More generally, any nonzero Laurent polynomial divisible by
\(v-v^{-1}\) fails \eqref{eq:purepointLaurentsign}:
it vanishes at \(v=-1\), whereas a nonzero effective
\((-1)^wf(-v)\) has positive value at \(v=1\).
For the pure realization
\eqref{eq:purepointLaurentconstruction}, this latter value
is exactly the total cohomology dimension.
Thus a raw nonzero signed/factored coefficient with a
positive exponent of \(v-v^{-1}\) needs a mixed or
normalized arithmetic comparison when \(v\) is assigned
weight one. Positivity of the stripped factor is a
different statement and is not contradicted.

\paragraph{A geometric mixed realization of the degeneration coefficient.}
The localization sequence for
\(\mathbb G_m\subset\mathbb A^1\)
\cite[Proposition~18.4(b), pp.~121--122]{Milne} gives
\[
 H_c^1((\mathbb G_m)_{\overline{\mathbb F}_q},
          \mathbb Q_\ell)=\mathbb Q_\ell,\qquad
 H_c^2((\mathbb G_m)_{\overline{\mathbb F}_q},
          \mathbb Q_\ell)=\mathbb Q_\ell(-1),
\]
with other groups zero. Indeed the closed point has
only \(H^0=\mathbb Q_\ell\), while affine-line compact
cohomology is \(\mathbb Q_\ell(-1)\) in degree two,
with the duality and Tate normalization in
\cite[Theorem~24.1(a),(b), p.~144]{Milne}.
The compact-support K\"unneth formula
\cite[Theorem~22.1 and Remark~22.3(b),
pp.~135--136]{Milne} therefore gives
\[
 K_0=R\Gamma_c((\mathbb G_m)^2_{\overline{\mathbb F}_q},
                   \mathbb Q_\ell(1)):
 \quad
 H^2K_0=\mathbb Q_\ell(1),\
 H^3K_0=\mathbb Q_\ell^{\,2},\
 H^4K_0=\mathbb Q_\ell(-1).
\]
Its Euler trace is
\[
                  q^{-n}-2+q^n
                  =q^{-n}(q^n-1)^2
                  =c_0(q^{n/2}).
\]
The coefficient sheaf \(\mathbb Q_\ell(1)\) is pure;
\(\mathbb Q_\ell(1)[2]\) is a pure weight-zero perverse
sheaf on the smooth two-dimensional torus.
Its compact pushforward to the point is mixed, as the
displayed cohomology shows. Thus the obstruction concerns
a claimed pure \emph{point complex}, not the existence
of ordinary geometric coefficient data.

\paragraph{A categorical corollary with state potentials.}
Fix finitely many states \(a\), integral potentials \(t_a\),
and generator weights \(w_i\in\mathbb Z\).
Suppose the \((a,b)\)-entry \(K_{i,ab}\) of each named
matrix kernel is a bounded pure point complex of weight
\begin{equation}\label{eq:purekernelstatepotential}
                            w_i+t_b-t_a.
\end{equation}
Compose kernels by tensor product and direct sum over
intermediate states. Tensor products of pure point
complexes add weights: in cohomological degree \(n\),
the K\"unneth summand \(H^jK\otimes H^{n-j}L\)
has weight \(w_K+j+w_L+n-j\).
A finite direct sum of objects of one weight preserves
that weight.
The entry of the word \(i_1\cdots i_L\) from \(a\) to
\(b\) is consequently pure of weight
\[
                        \sum_{j=1}^Lw_{i_j}+t_b-t_a.
\]
The potentials telescope along every path. All diagonal
entries, and hence their direct sum forming the closed-word
trace object, are pure of weight \(\sum_jw_{i_j}\).
Their Frobenius Euler traces are the usual matrix
convolution traces, by tensor multiplicativity and
direct-sum additivity.

The specified \fid{F57} word therefore excludes such a
kernel realization on \(N^\#\) if its closed-word Euler
trace is required to equal
\eqref{eq:rankfournegativeLaurent} under
\(v=q^{n/2}\).
This is an additional pure-word comparison, not a source
axiom. Purity of individual homogeneous summands of a
mixed entry is a weaker condition and is not excluded.
Likewise pushforward of a pure complex to the point along
a proper morphism of schemes of finite type over
\(\mathbb F_q\) cannot give this trace: BBD,
\S\S5.1.13--5.1.14, p.~128, preserves the upper weight
bound under \(f_!\) and the lower weight bound under
\(f_*\); properness identifies the two functors
\cite{BBD}.

The test applies to the scalar Euler trace of Frobenius
itself. It does not apply to an arbitrary inserted
operator trace. Neither the mixed Tate witness nor the
torus computation supplies the named nonstandard relation
maps, a designated canonical basis, or its specialization
filtration.

% Separate paragraph for Part III.
\paragraph{The pure-complex refinement in Grothendieck's language.}
The split-torus representation group separates effective
characters from Euler classes. SGA~3, I, 4.7.2--3,
and cote~90, scan~17 supply the action framework;
SGA~6, V, \S\S2--3, and cote~133, scan~18 supply
the coefficient representation and \(\lambda/\sigma\)
language \cite{SGA3,SGA6,Archive}.
No \(\lambda\)-ring structure on the noncommutative word
algebra is asserted.
The bounded coefficient complexes fit SGA~6,
I, Definition~4.7, pp.~121--122, and Example~4.8,
p.~122 \cite{SGA6}; their Weil actions and the
grading-to-weight comparison are extra specified data.
BBD, \S5.1.8, pp.~126--127, then imposes the
cohomological weight shift which yields
\eqref{eq:purepointLaurentsign}.
The nonproper torus example gives an ordinary mixed
compact pushforward. No archival positivity theorem or
designated nonstandard realization is inferred from it.
